% Capstone 4: quadratic sieve. Port of «Мини-курсовая - метод Квадратичного решета».
% The performance numbers come from the generated requirements.tex (make bench).
% The original's first reference is labelled "Discrete logarithm calculator" but links to
% alpertron's ECM.HTM (a copy-paste slip): it is cited as alpertron-ecm, "Integer
% factorization calculator".
\documentclass{ca-capstone}
\capstonerefs{alpertron-ecm,gmpy2,sagemath,hps}
\begin{document}
% Generated by `make bench CAPSTONE=quadratic-sieve` (taskgen.capstone report) from
% bench/results.json: reference time x slack on the reference machine.
% Do not edit: change bench/config.yaml and re-run the benchmark.
\newcommand{\ReqSlack}{10}
\newcommand{\ReqBits}{100}
\newcommand{\ReqSeconds}{10}

\capstonetitle{Факторизация методом Квадратичного решета}

\goal
Требуется реализовать разложение целого числа на множители методом \emph{Квадратичного
решета} "--- \emph{Quadratic sieve}, \emph{QS}.

\functionality
\begin{enumerate}
  \item Вычисление (1)~значения параметра~$B$, (2)~требуемого количества \emph{$B$-гладких}
    чисел и (3)~ожидаемой длины решета\footnote{параметр~$B$ и диапазон аргументов
    квадратичного полинома "--- длина просеиваемого отрезка~$l$ "--- зависят от раскладываемого
    числа~$N$ и выражаются в терминах функции~$L(N)$ (\emph{как?})} для заданного подлежащего
    разложению на множители числа~$N$
  \item Поиск \emph{$B$-гладких} чисел с помощью метода \emph{Квадратичного решета}
  \item Решение системы сравнений для отбора части \emph{$B$-гладких} квадратов, произведение
    которых даёт полный квадрат
  \item Нахождение нетривиального делителя заданного числа~$N$ с помощью НОД
\end{enumerate}

\requirements
Требования к реализации:
\begin{enumerate}
  \item корректность решения задачи факторизации целого числа
  \item приемлемая скорость работы (см. комментарий ниже)
\end{enumerate}

\textbf{Корректность работы} можно проверять, самостоятельно генерируя задачи, либо же
прибегая к помощи сторонних решателей, например, \cite{alpertron-ecm,sagemath}.

В данной работе важна не скорость, а корректность реализации и соответствие приведённым выше
требованиям. Однако тестовые примеры с числами размеров примерно
\textbf{\ReqBits~бит}\footnote{разумеется, трудоёмкость задачи факторизации конкретного числа
зависит от его структуры и сильно варьируется} должны выполняться за приемлемое время "---
не более \ReqSeconds~с.

\details
Так как цель работы "--- решение задачи факторизации целого числа, в своей реализации
\textbf{можно пользоваться библиотеками} для работы с большими целыми числами, а также пакетами
линейной алгебры.

\textbf{Квадратные корни} по модулю простого и степени простого числа~$p$ также можно находить
с помощью сторонних \textbf{библиотек}, либо же \textbf{перебором} (для небольших~$p$).

Рекомендуется использовать язык \emph{Python} и систему компьютерной алгебры
\emph{SageMath}~\cite{sagemath}.

Желательно сохранять промежуточные результаты решения задачи, такие как найденные на этапе
просеивания \emph{$B$-гладкие} квадраты, вид составленной системы сравнений и её решения,
результаты НОД отыскиваемых пробных делителей и факторизуемого числа~$N$ и проч.

Платформо-независимый, \textbf{переносимый код} писать \textbf{не требуется}, но во время сдачи
должна быть возможность проверить работу.

\schedule
Работа выполняется \textbf{индивидуально}. Срок выполнения "--- \textbf{месяц}. Работа должна
быть ваша: реализация должна быть \textbf{разработана и написана вами}, а не собрана из кусков
чужих работ.

Желательно хранить код в \emph{git}-репозитории на \emph{github}, \emph{bitbucket} или подобных
платформах. \textbf{Документация не требуется}.

\remarks
\begin{enumerate}
  \item Оценка значений параметров \emph{QS} приведена в разделе~3.7.1~\cite{hps}
  \item Просеивание нужно проводить только по простым числам (и их степеням) из фактор-базы,
    определяемой параметром~$B$
  \item Можно просеивать не по всем простым из фактор-базы
  \item Использование логарифмов и, как следствие, замена умножений вычитаниями, позволяет
    значительно ускорить поиск \emph{$B$-гладких} квадратов; недостатком такого подхода является
    необходимость работать с типом ограниченной точности (\emph{float}, \emph{double} и~т.\,п.)
  \item Можно использовать разные полиномы, не только $f(t) = t^2 - N$ (\emph{какие ещё?})
\end{enumerate}

\references
\end{document}
